\section{4AL Decomposition Algorithm}
\label{decomposition-algorithm}

\ifthenelse{\isundefined{\jfpda}}{
\begin{figure*}[t]
\begin{tabular}{ccc}
    \scalebox{0.7}{
      \ifthenelse{\isundefined{\jfpda}}{
\begin{minipage}{0.99\columnwidth}
}{
\begin{minipage}{0.49\columnwidth}
}
\vspace{-2.4cm}
\begin{algorithm}[H]
  \caption{Level 1 (Decomposing the Network)}
  \label{alg:level1}
%  \DontPrintSemicolon

  \KwIn{$LN$: Logical Network.}

  \KwOut{Approximation $V$ of expected value $V^*$ of attacking $LN$
    from controlled machine \start.}

  \acom{Decompose $LN$ into tree $DLN$ of biconnected components,
    rooted at \start; see text for ``clean-up''.} $DLN \gets
  $HopcroftTarjan$(LN)$\; 

  Set tree root to \start\ and \emph{clean-up} $LN$ and $DLN$\;

%  Set root of $DLN$ to \start\ and direct all arcs top-down in the
%  tree; Assign cut vertices to component closest to \start\;

  $\component_1, \dots, \component_k \gets$ a topological ordering of
  $DLN$\; 

  Intitialize pivoting reward $\pivotingrewardfn(N)$ for all $N \in
  LN$ to $0$\;

   \For{$i = k, ..., 1$}{ 

     \acom{Call Level 2 to attack each component.}  $V(\component_i)
     \gets$Level2$(\component_i,\pivotingrewardfn)$\;

     \acom{Propagate expected reward.} $N \gets$ the parent of
     $\component_i$ in $LN$\; 

     $\pivotingrewardfn(N) \gets \pivotingrewardfn(N) +
     V(\component_i)$\;

   }
   
  \Return{$\pivotingrewardfn(\start)$}
\end{algorithm}
\end{minipage}

    } &
    \scalebox{0.7}{
      \ifthenelse{\isundefined{\jfpda}}{
\begin{minipage}{0.99\columnwidth}
}{
\begin{minipage}{0.49\columnwidth}
}
\begin{algorithm}[H]
  \caption{Level 2 (Attacking Components)}
  \label{alg:level2}
%  \DontPrintSemicolon

  \KwIn{Biconnected component $\component$, reward function
    $\pivotingrewardfn$.}

  \KwOut{Approximation $V$ of expected value $V^*$ of attacking
    $\component$, given its parent is controlled and its pivoting
    rewards are $\pivotingrewardfn$.}

%  \acom{Consider all simple paths -- with no repeated vertices --
%    between entry vertices and (pivoting) reward vertices in
%    $\component$.} 
  
  $\reward \gets 0$\;

  \acom{Account for each rewarded vertex $N$.}

  \While{$ \exists N \in \component$ s.t.\ $\rewardfn(N) > 0$ or
    $\pivotingrewardfn(N)>0$}{

    $P \gets \langle \rangle$; $\reward(P) \gets 0$; $P(N) \gets P$\;

    \acom{Maximize over all simple paths (no repeated vertices) from
      an entry vertex to $N$.}

    \ForEach{simple path $P$ of the form $\xrightarrow{F_0} N_1
      \xrightarrow{F_1} N_2 \dots \xrightarrow{F_{k-1}} N_k = N$ where
      $N_1, \dots, N_k \in \component$ and $N_1 \in
      \component_\start$}{

      \acom{Propagate rewards along $P$, calling Level 3 for attack on
        each subnetwork.}

      $\reward(P) \gets 0$\; 

      \For{$i = k, ..., 1$}{

        $\reward(P) \gets$ Level3$(N_{i},F_{i-1},\pivotingrewardfn(N_{i}),\reward(P))$\; 

      }
      
      $P(N) \gets \text{arg max}(\reward(P(N)),\reward(P))$\;

    }

    $\reward \gets \reward + \reward(P(N))$\;

    $\rewardfn(N_i), \pivotingrewardfn(N_i) \gets 0$ for all vertices $N_i$ on $P(N)$\;

  }

  \Return{$\reward$}

\end{algorithm}
\end{minipage}

    } &
    \scalebox{0.7}{
      \ifthenelse{\isundefined{\jfpda}}{
\begin{minipage}{0.99\columnwidth}
}{
\begin{minipage}{0.49\columnwidth}
}
\vspace{-2.43cm}
\begin{algorithm}[H]
  \caption{Level 3 (Attacking Subnetworks)}
  \label{alg:level3}
%  \DontPrintSemicolon

  \KwIn{Firewall $F$, subnetwork $N$, rewards $\pivotingreward$,
    $\pathreward$.}

  \KwOut{Approximation $V$ of expected value $V^*$ of attacking $N$
    through $F$, given $F$ is reached, $N$'s pivoting reward is
    $\pivotingreward$, and the path reward behind $N$ is
    $\pathreward$.}

  $\reward \gets 0$\;

  \acom{Maximize over reward obtained when hacking first into a
    particular machine $m \in N$.}

  \ForEach{$m \in N$}{

    $\reward(m) \gets \rewardfn(m)$\;

    \acom{After breaking $m$, we can pivot behind $N$, and reach all
      $m \neq m' \in N$ without $F$.}

    $\reward(m) \gets \reward(m) + \pivotingreward + \pathreward$\;

    \ForEach{$m \neq m' \in N$}{

      $\reward(m) \gets \reward(m) +
      $Level4$(m',\emptyfirewall,\rewardfn(m'))$\;

    }

    $\reward \gets \max(\reward,$ Level4$(m,F,\reward(m)))$\;

  }

  \Return{$\reward$}
\end{algorithm}
\end{minipage}

    } \\[0.1cm]
    (a) $LN$ as tree of components $C$. & 
    (b) Paths for attacking $C_1$. & 
    (c) Attacking $N_3$ from $N_1$, using $m$ first.   
  \end{tabular}
  \vspace{-0.35cm}
  \caption{Illustration of Levels 1, 2, and 3 (from left to right) of
    the 4AL algorithm.}
  \label{fig:4ALlevels123}
  \vspace{-0.5cm}
\end{figure*}
}{
  \begin{figure}[t]
    \begin{minipage}{0.45\linewidth}
      \centering
      \scalebox{0.7}{
        \ifthenelse{\isundefined{\jfpda}}{
\begin{minipage}{0.99\columnwidth}
}{
\begin{minipage}{0.49\columnwidth}
}
\vspace{-2.4cm}
\begin{algorithm}[H]
  \caption{Level 1 (Decomposing the Network)}
  \label{alg:level1}
%  \DontPrintSemicolon

  \KwIn{$LN$: Logical Network.}

  \KwOut{Approximation $V$ of expected value $V^*$ of attacking $LN$
    from controlled machine \start.}

  \acom{Decompose $LN$ into tree $DLN$ of biconnected components,
    rooted at \start; see text for ``clean-up''.} $DLN \gets
  $HopcroftTarjan$(LN)$\; 

  Set tree root to \start\ and \emph{clean-up} $LN$ and $DLN$\;

%  Set root of $DLN$ to \start\ and direct all arcs top-down in the
%  tree; Assign cut vertices to component closest to \start\;

  $\component_1, \dots, \component_k \gets$ a topological ordering of
  $DLN$\; 

  Intitialize pivoting reward $\pivotingrewardfn(N)$ for all $N \in
  LN$ to $0$\;

   \For{$i = k, ..., 1$}{ 

     \acom{Call Level 2 to attack each component.}  $V(\component_i)
     \gets$Level2$(\component_i,\pivotingrewardfn)$\;

     \acom{Propagate expected reward.} $N \gets$ the parent of
     $\component_i$ in $LN$\; 

     $\pivotingrewardfn(N) \gets \pivotingrewardfn(N) +
     V(\component_i)$\;

   }
   
  \Return{$\pivotingrewardfn(\start)$}
\end{algorithm}
\end{minipage}

      }
      \\
      (a) $LN$ as tree of components $C$.
    \end{minipage}
    \hfill  
    \begin{minipage}{0.45\linewidth}
      \centering
      \scalebox{0.7}{
        \ifthenelse{\isundefined{\jfpda}}{
\begin{minipage}{0.99\columnwidth}
}{
\begin{minipage}{0.49\columnwidth}
}
\begin{algorithm}[H]
  \caption{Level 2 (Attacking Components)}
  \label{alg:level2}
%  \DontPrintSemicolon

  \KwIn{Biconnected component $\component$, reward function
    $\pivotingrewardfn$.}

  \KwOut{Approximation $V$ of expected value $V^*$ of attacking
    $\component$, given its parent is controlled and its pivoting
    rewards are $\pivotingrewardfn$.}

%  \acom{Consider all simple paths -- with no repeated vertices --
%    between entry vertices and (pivoting) reward vertices in
%    $\component$.} 
  
  $\reward \gets 0$\;

  \acom{Account for each rewarded vertex $N$.}

  \While{$ \exists N \in \component$ s.t.\ $\rewardfn(N) > 0$ or
    $\pivotingrewardfn(N)>0$}{

    $P \gets \langle \rangle$; $\reward(P) \gets 0$; $P(N) \gets P$\;

    \acom{Maximize over all simple paths (no repeated vertices) from
      an entry vertex to $N$.}

    \ForEach{simple path $P$ of the form $\xrightarrow{F_0} N_1
      \xrightarrow{F_1} N_2 \dots \xrightarrow{F_{k-1}} N_k = N$ where
      $N_1, \dots, N_k \in \component$ and $N_1 \in
      \component_\start$}{

      \acom{Propagate rewards along $P$, calling Level 3 for attack on
        each subnetwork.}

      $\reward(P) \gets 0$\; 

      \For{$i = k, ..., 1$}{

        $\reward(P) \gets$ Level3$(N_{i},F_{i-1},\pivotingrewardfn(N_{i}),\reward(P))$\; 

      }
      
      $P(N) \gets \text{arg max}(\reward(P(N)),\reward(P))$\;

    }

    $\reward \gets \reward + \reward(P(N))$\;

    $\rewardfn(N_i), \pivotingrewardfn(N_i) \gets 0$ for all vertices $N_i$ on $P(N)$\;

  }

  \Return{$\reward$}

\end{algorithm}
\end{minipage}

      }
      \\
      (b) Paths for attacking $C_1$.
    \end{minipage}

    \medskip
    
    \begin{minipage}{1\linewidth}
      \centering
      \scalebox{0.7}{
        \ifthenelse{\isundefined{\jfpda}}{
\begin{minipage}{0.99\columnwidth}
}{
\begin{minipage}{0.49\columnwidth}
}
\vspace{-2.43cm}
\begin{algorithm}[H]
  \caption{Level 3 (Attacking Subnetworks)}
  \label{alg:level3}
%  \DontPrintSemicolon

  \KwIn{Firewall $F$, subnetwork $N$, rewards $\pivotingreward$,
    $\pathreward$.}

  \KwOut{Approximation $V$ of expected value $V^*$ of attacking $N$
    through $F$, given $F$ is reached, $N$'s pivoting reward is
    $\pivotingreward$, and the path reward behind $N$ is
    $\pathreward$.}

  $\reward \gets 0$\;

  \acom{Maximize over reward obtained when hacking first into a
    particular machine $m \in N$.}

  \ForEach{$m \in N$}{

    $\reward(m) \gets \rewardfn(m)$\;

    \acom{After breaking $m$, we can pivot behind $N$, and reach all
      $m \neq m' \in N$ without $F$.}

    $\reward(m) \gets \reward(m) + \pivotingreward + \pathreward$\;

    \ForEach{$m \neq m' \in N$}{

      $\reward(m) \gets \reward(m) +
      $Level4$(m',\emptyfirewall,\rewardfn(m'))$\;

    }

    $\reward \gets \max(\reward,$ Level4$(m,F,\reward(m)))$\;

  }

  \Return{$\reward$}
\end{algorithm}
\end{minipage}

      }
      \\
      (c) Attacking $N_3$ from $N_1$, using $m$ first.   
    \end{minipage}
  %\vspace{-0.35cm}
  \caption{Illustration of Levels 1, 2, and 3 (from left to right) of
    the 4AL algorithm.}
  \label{fig:4ALlevels123}
  %\vspace{-0.5cm}
\end{figure}
}

\begin{figure*}[t]
\centering
\begin{tabular}{cc}
%\vspace{-0.1cm}
\ifthenelse{\isundefined{\jfpda}}{
\begin{minipage}{0.99\columnwidth}
}{
\begin{minipage}{0.49\columnwidth}
}
\vspace{-2.4cm}
\begin{algorithm}[H]
  \caption{Level 1 (Decomposing the Network)}
  \label{alg:level1}
%  \DontPrintSemicolon

  \KwIn{$LN$: Logical Network.}

  \KwOut{Approximation $V$ of expected value $V^*$ of attacking $LN$
    from controlled machine \start.}

  \acom{Decompose $LN$ into tree $DLN$ of biconnected components,
    rooted at \start; see text for ``clean-up''.} $DLN \gets
  $HopcroftTarjan$(LN)$\; 

  Set tree root to \start\ and \emph{clean-up} $LN$ and $DLN$\;

%  Set root of $DLN$ to \start\ and direct all arcs top-down in the
%  tree; Assign cut vertices to component closest to \start\;

  $\component_1, \dots, \component_k \gets$ a topological ordering of
  $DLN$\; 

  Intitialize pivoting reward $\pivotingrewardfn(N)$ for all $N \in
  LN$ to $0$\;

   \For{$i = k, ..., 1$}{ 

     \acom{Call Level 2 to attack each component.}  $V(\component_i)
     \gets$Level2$(\component_i,\pivotingrewardfn)$\;

     \acom{Propagate expected reward.} $N \gets$ the parent of
     $\component_i$ in $LN$\; 

     $\pivotingrewardfn(N) \gets \pivotingrewardfn(N) +
     V(\component_i)$\;

   }
   
  \Return{$\pivotingrewardfn(\start)$}
\end{algorithm}
\end{minipage}

&
\ifthenelse{\isundefined{\jfpda}}{
\begin{minipage}{0.99\columnwidth}
}{
\begin{minipage}{0.49\columnwidth}
}
\begin{algorithm}[H]
  \caption{Level 2 (Attacking Components)}
  \label{alg:level2}
%  \DontPrintSemicolon

  \KwIn{Biconnected component $\component$, reward function
    $\pivotingrewardfn$.}

  \KwOut{Approximation $V$ of expected value $V^*$ of attacking
    $\component$, given its parent is controlled and its pivoting
    rewards are $\pivotingrewardfn$.}

%  \acom{Consider all simple paths -- with no repeated vertices --
%    between entry vertices and (pivoting) reward vertices in
%    $\component$.} 
  
  $\reward \gets 0$\;

  \acom{Account for each rewarded vertex $N$.}

  \While{$ \exists N \in \component$ s.t.\ $\rewardfn(N) > 0$ or
    $\pivotingrewardfn(N)>0$}{

    $P \gets \langle \rangle$; $\reward(P) \gets 0$; $P(N) \gets P$\;

    \acom{Maximize over all simple paths (no repeated vertices) from
      an entry vertex to $N$.}

    \ForEach{simple path $P$ of the form $\xrightarrow{F_0} N_1
      \xrightarrow{F_1} N_2 \dots \xrightarrow{F_{k-1}} N_k = N$ where
      $N_1, \dots, N_k \in \component$ and $N_1 \in
      \component_\start$}{

      \acom{Propagate rewards along $P$, calling Level 3 for attack on
        each subnetwork.}

      $\reward(P) \gets 0$\; 

      \For{$i = k, ..., 1$}{

        $\reward(P) \gets$ Level3$(N_{i},F_{i-1},\pivotingrewardfn(N_{i}),\reward(P))$\; 

      }
      
      $P(N) \gets \text{arg max}(\reward(P(N)),\reward(P))$\;

    }

    $\reward \gets \reward + \reward(P(N))$\;

    $\rewardfn(N_i), \pivotingrewardfn(N_i) \gets 0$ for all vertices $N_i$ on $P(N)$\;

  }

  \Return{$\reward$}

\end{algorithm}
\end{minipage}

\\
\ifthenelse{\isundefined{\jfpda}}{
\begin{minipage}{0.99\columnwidth}
}{
\begin{minipage}{0.49\columnwidth}
}
\vspace{-2.43cm}
\begin{algorithm}[H]
  \caption{Level 3 (Attacking Subnetworks)}
  \label{alg:level3}
%  \DontPrintSemicolon

  \KwIn{Firewall $F$, subnetwork $N$, rewards $\pivotingreward$,
    $\pathreward$.}

  \KwOut{Approximation $V$ of expected value $V^*$ of attacking $N$
    through $F$, given $F$ is reached, $N$'s pivoting reward is
    $\pivotingreward$, and the path reward behind $N$ is
    $\pathreward$.}

  $\reward \gets 0$\;

  \acom{Maximize over reward obtained when hacking first into a
    particular machine $m \in N$.}

  \ForEach{$m \in N$}{

    $\reward(m) \gets \rewardfn(m)$\;

    \acom{After breaking $m$, we can pivot behind $N$, and reach all
      $m \neq m' \in N$ without $F$.}

    $\reward(m) \gets \reward(m) + \pivotingreward + \pathreward$\;

    \ForEach{$m \neq m' \in N$}{

      $\reward(m) \gets \reward(m) +
      $Level4$(m',\emptyfirewall,\rewardfn(m'))$\;

    }

    $\reward \gets \max(\reward,$ Level4$(m,F,\reward(m)))$\;

  }

  \Return{$\reward$}
\end{algorithm}
\end{minipage}

&
\ifthenelse{\isundefined{\jfpda}}{
\begin{minipage}{0.99\columnwidth}
}{
\begin{minipage}{0.49\columnwidth}
}
\vspace{0.4cm}
\begin{algorithm}[H]
  \caption{Level 4 (Attacking Individual Machines)}
  \label{alg:level4}
%  \DontPrintSemicolon

  \KwIn{Firewall $F$, machine $m$, reward $\reward$.}

  \KwOut{Approximation $V$ of expected value $V^*$ of attacking $m$
    through $F$, given $m$ is reached and the current reward of
    breaking it is $\reward$.}

  \If{$(m,F,\reward)$ is cached}{ \Return{$V(m,F,\reward)$} }

%  \acom{Create POMDP as usual and run.}  
  $M \gets$createPOMDP$(m,F,\reward)$\; 

  $V \gets$solvePOMDP$(M)$\;

%  \acom{Remember this solution} 
  Cache $(m,F,\reward)$ with $V$\;

  \Return{$V$}

\end{algorithm}
\end{minipage}

\end{tabular}
\vspace{-0.2cm}\caption{4AL algorithm, pseudo-code.}
\label{fig:4ALcode}
\vspace{-0.5cm}
\end{figure*}

As hinted, POMDPs do not scale to large networks (cf.\ the experiments
in the next section). We now present an approach using decomposition
and approximation to overcome this problem, relying on POMDPs only to
attack individual machines. The approach is called \emph{4AL} since it
addresses network attack at 4 different levels of abstraction. 4AL is
a POMDP solver specialized to attack planning as addressed here. Its
input are the logical network $LN$ and POMDP models encoding attacks
on individual machines. Its output is a policy (an attack) for the
global POMDP encoding $LN$, as well as an approximation of the value
of the global value function. We next overview the algorithm, then
fill in some technical details. To simplify the presentation, we will
focus on the approximation of the value function, and outline only
briefly how to construct the policy.




\subsection{4AL Overview and Basic Properties}
\label{decomposition-algorithm:overview}


The four levels of 4AL are: (1) \emph{Decomposing the Network}, (2)
\emph{Attacking Components}, (3) \emph{Attacking Subnetworks}, and (4)
\emph{Attacking Individual Machines}. We outline these levels in
turn before providing technical details. Figure~\ref{fig:4ALlevels123} provides illustrations.

\vspace{-0.05cm}
\begin{itemize}
\item \textbf{Level~1:} Decompose the logical network $LN$ into a tree
  of biconnected components, rooted at \start. In reverse topological
  order, call Level~2 on each component; propagate the outcomes
  upwards in the tree.
\end{itemize}
\vspace{-0.05cm}

Every graph decomposes into a unique tree of biconnected components
\cite{HopTar-cACM73}. A biconnected component is a sub-graph that
remains connected when removing any one vertex. In pentesting,
intuitively this means that there is more than one possibility (more
than one path) to attack the subnetworks within the component,
requiring to reason about the component as a whole (which is the job
of Level~2). By contrast, if removing subnetwork $X$ (\eg, $N_2$
in Figure~\ref{fig:4ALlevels123} (b)) makes the graph fall apart into
two separate sub-graphs ($C_2$ vs.\ the rest of $LN$, compare also
Figure~\ref{fig:4ALlevels123} (a)), then \emph{all} attacks from
\start\ to one of these sub-graphs ($C_2$ here) must first traverse
$X$ ($N_2$ here). Thus the overall expected value of the attack can be
computed by (1) computing the value of attacking that sub-graph
($C_2$) alone, %
%given $X$ is already under control, 
and (2) adding the result as a \emph{pivoting reward} to the reward of
breaking into $X$ ($N_2$). In other words, we ``propagate the outcomes
upwards'' in the tree displayed in Figure~\ref{fig:4ALlevels123}
(a). 

It is important to note that this tree decomposition will typically
result in a huge reduction of complexity. Biconnected components in
$LN$ arise only from clusters of more than $2$ subnetworks sharing a
common (physical) firewall machine. Such clusters tend to be small. In
the real-world test scenario used by Core Security and in our
experiment here, there is only one cluster, of size $3$. In case there
are no clusters at all, $LN$ is a tree and 4AL Level~2 trivializes
completely.

\vspace{-0.05cm}
\begin{itemize}
\item \textbf{Level~2:} Given component $\component$, consider, for
  each rewarded subnetwork $N \in \component$, all paths $P$ in
  $\component$ that reach $N$. Backwards along each $P$, call Level~3
  on each subnetwork and associated firewall. Choose the best path for
  each $N$. Aggregate these path values over all $N$, by summing up
  but disregarding rewards that were already accounted for by a
  previous path in the sum.
\end{itemize}
\vspace{-0.05cm}

In case a biconnected component $\component$ contains more than one
subnetwork, to obtain the best attack on $\component$, in general we
have no choice but to encode the entire component as a POMDP. Since
that is not feasible, Level~2 considers individual ``attack paths''
within $\component$. Any single path $P$ is equivalent to a sequence
of attacks on individual subnetworks; these attacks are evaluated
using Level~3. We consider the rewarded vertices $N$ in separation,
enumerating the attack paths and choosing a best one.  The values of
the best paths are aggregated over all $N$ in a conservative
(pessimistic) manner, by accounting for each reward at most once. A
strict under-estimation occurs in case the best paths for some
rewarded vertices are not disjoint: then these attacks share some of
their cost, so a combined attack has a higher expected reward than the
sum of independent attacks.

In Figure~\ref{fig:4ALlevels123} (b), $N_2$ and $N_3$ have a pivoting
reward because they allow to reach the components $C_2$ and $C_3$
respectively. If the best paths for both $N_2$ and $N_3$ go via $N_1$
(because the firewall $F^*_3$ is very strict), then these paths are
not disjoint, duplicating the effort for breaking into $N_1$.

Obviously, enumerating attack paths within $\component$ is exponential
in the size of $\component$. This is the only point in 4AL---apart of
course from calls to the POMDP solver---that has worst-case
exponential runtime. In practice, biconnected components are typically
small, \cf\ the above.

\vspace{-0.05cm}
\begin{itemize}
\item \textbf{Level~3:} Given a subnetwork $N$ and a firewall $F$
  through which to attack $N$, for each machine $m \in N$ approximate
  the reward obtained when attacking $m$ first. For this, modify $m$'s
  reward to take into account that, after breaking $m$, we are behind
  $F$: call Level~4 to obtain the values of all $m' \neq m$ with an
  empty firewall; then add these values, plus any pivoting reward, to
  the reward of $m$ and call Level~4 on this modified $m$ with
  firewall $F$. Maximize the resulting value over all $m \in N$.
\end{itemize}
\vspace{-0.05cm}

Consider Figure~\ref{fig:4ALlevels123} (c). When attacking $N$ (here,
$N_3$) from some machine behind the firewall $F$ (here, $F^1_3$), we
have to choose which machine inside $N$ to attack. Given we commit to
one such choice $m$, the attack problem becomes that of breaking into
$m$ and afterwards exploiting the direct connection to any $m \neq m'
\in N$, and any descendant network (here, $C_3$) we can now pivot
to. As described, that can be dealt with by combining attacks on
individual machines with modified rewards. (The pivoting reward for
descendant networks is computed beforehand by Levels 1 and 2.)


Like Level~2, Level~3 makes a conservative approximation. It fixes a
choice of which $m \in N$ to attack. By contrast, the best strategy
may be to switch between different $m \in N$ depending on the success
of the attack so far. For example, if one exploit is very likely to
succeed, then it may pay off to try this on all $m$ first, before
trying anything else.


\vspace{-0.05cm}
\begin{itemize}
\item \textbf{Level~4:} Given a machine $m$ and a firewall $F$, model
  the single-machine attack planning problem as a POMDP, and run an
  off-the-shelf POMDP solver. Cache known results to avoid duplicate
  effort.
\end{itemize}
\vspace{-0.05cm}

This last step should be self-explanatory. The POMDP model is created
as described earlier. Note that Level~3 may, during the execution of
4AL, call the same machine with the same firewall more than once. For
example, in Figure~\ref{fig:4ALlevels123} (c), when we switch to
attacking $m'_1$ instead of $m$, the call of Level~4 with $m'_k$ and
an empty firewall is repeated.

Summing up, 4AL has low-order polynomial runtime except for the
enumeration of paths within biconnected components (Level~2), and
solving single-machine POMDPs (Level~4). The decomposition at Level~1
incurs no information loss. Levels 2 and 3 make conservative
approximations, so, if the POMDP solutions are conservative (\eg,
optimal), then the overall outcome of 4AL is conservative as well. 








\subsection{Technicalities}
\label{decomposition-algorithm:technicalities}




To provide a more detailed understanding of 4AL, we now discuss
pseudo-code for the algorithm, provided in
Figure~\ref{fig:4ALcode}. Consider first
Algorithm~\ref{alg:level1}. It should be clear how the overall
structure of the algorithm corresponds to our previous discussion. It
calls the linear-time algorithm by Hopcroft and Tarjan
\shortcite{HopTar-cACM73} (hereafter, HT) to find the
decomposition. The loop $i=k,\dots,1$ processes the components in
reverse topological order. The pivoting reward function
$\pivotingrewardfn$ encodes the propagation of rewards upwards in the
tree; this should be self-explanatory apart for the expression ``the
parent'' of $C_i$ in $LN$. The latter relies on the fact that, after
``clean-up'' (line 2), each component has exactly one such parent.

To explain the clean-up, note first that HT works on undirected
graphs; when applying it, we ignore the direction of the arcs in
$LN$. The outcome is an undirected tree of biconnected components,
where the \emph{cut vertices}---those vertices removing which makes
the graph break apart---are shared between several components. In
Figure~\ref{fig:4ALlevels123} (b), \eg, $N_2$ prior to the clean-up
belongs to both, $C_1$ and $C_2$. The clean-up sets the root of the
tree to \start, and assigns each cut-vertex to the component closest
to \start\ (\eg, $N_2$ is assigned to $C_1$); \start\ itself is turned
into a separate component. Re-introducing the direction of arcs in
$LN$, we then prune vertices not reachable from \start. Next, we
remove arcs that cannot participate in any non-redundant attack path
starting in \start. Since moving \emph{towards} \start\ in the
decomposition tree necessarily leads any attack back to a vertex it
has visited (broken into) already, after such removal the arcs between
components form a directed tree as in Figure~\ref{fig:4ALlevels123}
(a). Each non-root component $\component_i$ (\eg, $C_3$) has exactly
one parent component $\component$ in the cleaned-up tree (\eg,
$C_1$). The respective subnetwork $N \in \component$ (\eg, $N_3$) is a
cut vertex in $LN$. Thus, as claimed above, $N$ is the \emph{only}
vertex, in $LN$, that connects into $\component_i$.

Obviously, all attacks on $\component_i$ must pass through its parent
$N$. Further, the vertices and arcs removed by clean-up cannot be part
of an optimal attack. Thus Level~1 is loss-free. To state this---and
the other properties of 4AL---formally, we need some notations. We
will use $V^*$ to denote the real (optimal) expected value of an
attack, and $V$ to denote the 4AL approximation. The attacked object
is given as the argument. For example, $V^*(LN)$ is the expected value
of attacking $LN$; $V(C,\pivotingrewardfn)$ is the outcome of running
4AL Level~2 on component $\component$ and pivoting reward function
$\pivotingrewardfn$.

\begin{proposition}\label{pro:level1correct} 
Let $LN$ be a logical network. Say that, for all calls to 4AL Level~2
made by 4AL Level~1 when run on $LN$, we have $V(C,\pivotingrewardfn)
= V^*(C,\pivotingrewardfn)$. Then $V(LN) = V^*(LN)$. If
$V(C,\pivotingrewardfn) \leq V^*(C,\pivotingrewardfn)$ for all calls
to 4AL Level~2, then $V(LN) \leq V^*(LN)$.
\end{proposition}


Consider now Algorithm~\ref{alg:level2}. Our previous description was
imprecise in omitting the additional algorithm argument
$\pivotingrewardfn$. This integrates with the algorithm by being
passed on, for every subnetwork on the paths we consider (line 7), to
Algorithm~\ref{alg:level3} which adds it to the reward obtained for
hacking into that subnetwork (Algorithm~\ref{alg:level3} line 4).

$\reward$ aggregates the values (lines 1, 9), over all rewarded
subnetworks $N$. This aggregation is made conservative by removing all
rewards---pivoting rewards as well as the own rewards of the
individual machines involved---that have already been accounted for
(line 10). Regarding the individual machines,
Algorithm~\ref{alg:level2} uses the shorthands (a) $\rewardfn(N) > 0$
(line 2) and (b) $\rewardfn(N) \leftarrow 0$ (line 10); (a) means that
there exists $m \in N$ so that $\rewardfn(m) > 0$; (b) means that
$\rewardfn(m) \leftarrow 0$ for all $m \in N$. Regarding pivoting
rewards, note that line 10 of Algorithm~\ref{alg:level2} modifies the
function $\pivotingrewardfn$ maintained by
Algorithm~\ref{alg:level1}. This does not lead to conflicts because,
at the time when Algorithm~\ref{alg:level1} calls
Algorithm~\ref{alg:level2} on component $C$, all descendants of $C$ in
$LN$ have already been processed, and thus in particular
Algorithm~\ref{alg:level1} will make no further updates to the value
of $\pivotingrewardfn(N)$, for any $N \in C$.

By $\component_\start$ (line 4) we denote the set $\{N \in \component
\mid \exists N' \in LN, N' \not \in \component: (N',N) \in LN\}$ of
subnetworks that serve as an entry into $\component$ (\eg, $N_1$ and
$N_3$ for $C_1$ in Figure~\ref{fig:4ALlevels123} (b)).  Note in line 4
that the path $P$ starts with a firewall $F_0$. To understand this,
consider the situation addressed. The algorithm assumes that the
parent $N$ of $\component$ (\start, for component $C_1$ in
Figure~\ref{fig:4ALlevels123} (b)) is under control. But then, to
break into $\component$, we still need to traverse an arc from $N$
into $\component$. $F_0$ is the firewall on the arc chosen by $P$
($F^*_1$ or $F^*_3$ in Figure~\ref{fig:4ALlevels123} (b)).

The calls to Level~3 (line 7) comprise the network $N_i$ to be hacked
into, the firewall $F_{i-1}$ that must be traversed for doing so, the
pivoting reward of $N_i$, as well as the ongoing path reward
$\reward(P)$ which gets propagated backwards along the path. Clearly,
this is equivalent to the sequence of attacks required to execute $P$,
and harvesting all pivoting rewards associated with such an
attack. Thus, with the conservativeness of the aggregation across the
subnetworks $N$, we get:

\begin{proposition}\label{pro:level2conservative} 
Let $\component$ be a biconnected component, and let
$\pivotingrewardfn$ be a pivoting reward function. Say that, for all
calls to 4AL Level~3 made by 4AL Level~2 when run on
$(\component,\pivotingrewardfn)$, we have
$V(F,N,\pivotingreward,\pathreward) \leq
V^*(F,N,\pivotingreward,\pathreward)$. Then $V(C,\pivotingrewardfn)
\leq V^*(C,\pivotingrewardfn)$.
\end{proposition}


Algorithms~\ref{alg:level3} and~\ref{alg:level4} should be
self-explanatory, given our previous discussion. Just note that the
pivoting reward $\pivotingreward$ is represented by the arc from $m$
to $C_3$ in Figure~\ref{fig:4ALlevels123} (c), which is accounted for
by simply adding it to the value of $m$ (Algorithm~\ref{alg:level3}
line 4). The path reward $\pathreward$ (not illustrated in
Figure~\ref{fig:4ALlevels123} (c)) is also added to the value of $m$
(Algorithm~\ref{alg:level3} line 4). Max'ing over attacks on the
individual machines $m$ is, obviously, a conservative approximation
because attack strategies are free to choose $m$. Thus:

\begin{proposition}\label{pro:level3conservative} 
Let $F$ be a firewall, let $N$ be a subnetwork, let $\pivotingreward$
be a pivoting reward, and let $\pathreward$ be a path reward. Say
that, for all calls to 4AL Level~4 made by 4AL Level~3 when run on
$(F,N,\pivotingreward,\pathreward)$, we have $V(F,m,\reward) \leq
V^*(F,m,\reward)$. Then $V(F,N,\pivotingreward,\pathreward) \leq
V^*(F,N,\pivotingreward,\pathreward)$.
\end{proposition}











\subsection{Policy Construction} 


At Level~1, the global policy is constructed from the Level~2 policies
simply by following the tree decomposition: starting at the tree root,
we execute the Level~2 policies for all reached components (in any
order); once a hack into a component succeeds, the respective children
components become reached. At Level~2, i.e., within a bi-connected
component $C$, the policy corresponds to the set of paths $P$
considered by Algorithm~\ref{alg:level2}. Each $P$ is processed in
turn. For each node $N$ in $P$ (until failure to enter that
subnetwork), we call the corresponding Level~3 policy.

At Level~3, i.e., considering a single subnetwork $N$, our policy
simply attacks the machine $m \in N$ that yielded the maximum in
Algorithm~\ref{alg:level3}. The policy first attacks $m$ through the
firewall, using the respective Level~4 policy. In case the attack
succeeds, the policy attacks the remaining machines $m' \in N$ in any
order (i.e., for each $m'$, we perform the associated Level~4 policy
until termination). At Level~4, the policy is the POMDP policy
returned by our POMDP solver.



